%% =====================================================================
%%  B1-T-05-01 -- what B1 contributes to "The Real Objection"
%%  -------------------------------------------------------------------
%%  LAYER A: APPEND-ONLY. Written once, then left alone. Material found
%%  only here sits in the `unique` slot and is protected by rule R10.
%% =====================================================================
%% @kind: source
%% @id: B1-T-05-01
%% @book: B1
%% @topic: T-05-01
%% @locator: Tactic 6, pp. 71--76
%% @status: distilled
%% @unique: yes
%% @integrated: yes
%% @updated: 2026-09-19

\dossier{B1}{T-05-01}{\bookshort{B1}}{Tactic 6, pp. 71--76}

\begin{distilled}
  \item Picking a person apart to discover real motives and feelings is presented as the prerequisite of influencing them at all.
  \item The questions the practitioner must answer are stated plainly: why will he not sign, why will she not do what is asked.
  \item Without the real objection you do not know what obstacle stands in the way, and effort spent on the stated objection is described as swinging in the dark --- most blows fall wide of the target.
  \item Finding the true objection is said to give you something you can actually work with.
\end{distilled}

\begin{unique}
  \item The operational distinction between the stated and the real objection, and the diagnostic value of the stated one: effort directed at it misses, and the pattern of misses identifies the real target.
\end{unique}
